\section{Event-level QA}
\label{sec:evtqa}
All figures in Secs.~\ref{sec:evtqa}--\ref{sec:tofqa} are produced by
\fpath{analysis_note/scripts/make_figures.C} from \fpath{data/qa_test_10k.root} (copy of
\fpath{output/qa_test_10k.root} on the node). Unless stated otherwise they show the
\nEvtSel{} selected events; the event count printed above each frame is the number of selected events.
Errors are statistical.

\begin{table}[h]
\centering
\caption{Event cut flow for the test sample (\fcmd{hEvtCut}; generated table \fpath{tables/evtcut.tex}).}
\label{tab:evtcut}
\input{tables/evtcut.tex}
\end{table}

\begin{table}[h]
\centering
\caption{Trigger IDs stored in the picoDst for all \nEvtAll{} events read (an event can carry several).
IDs with fewer than 50 events are summed. \nTrigIds{} distinct IDs occur in total
(full list in \fpath{logs/numbers.txt}).}
\label{tab:trigger}
\input{tables/trigger.tex}
\end{table}

\paragraph{Trigger.} \nEvtTrig{} of \nEvtAll{} events carry 640051, the only minimum-bias ID present.
IDs 6, 35, 36 and 37 also appear in a large fraction of events.
\caveat{The meaning of trigger IDs 6, 35, 36, 37, 640601 and 6408xx was not looked up. They do not
affect the selection, which only requires one of the minimum-bias IDs.}

\paragraph{Vertex.} The $V_z$ distribution before cuts is broad (RMS \rmsVzAll\,cm, mean \meanVzAll\,cm),
so the $|V_z|<70$\,cm cut is the largest loss (Table~\ref{tab:evtcut}, Fig.~\ref{fig:evt_vz}).
The transverse vertex is at $(\langle V_x\rangle,\langle V_y\rangle)=(\meanVx,\meanVy)$\,cm with RMS
\rmsVx{} and \rmsVy\,cm (Fig.~\ref{fig:evt_vxvy}), well inside the $V_r<2$\,cm circle centred at the origin.

\paragraph{VPD vertex.} For \fracVpdValid\% of the selected events $V_z^{\rm VPD}$ lies within $\pm200$\,cm;
the remainder has no valid VPD vertex. These events, together with events with
$|\Delta V_z|>20$\,cm, populate the under-/overflow of Fig.~\ref{fig:evt_vzdiff} (\nVzDiffUnder{} / \nVzDiffOver{} events). A Gaussian fit to the core ($|\Delta V_z|<3$\,cm) of
$\Delta V_z = V_z-V_z^{\rm VPD}$ gives mean \vzDiffMean\,cm and $\sigma=\vzDiffSigma$\,cm;
only \fracVzDiffThree\% of the selected events have $|\Delta V_z|<3$\,cm.

\begin{figure}[H]
\centering
\figpanel{0.48}{evt_vz}{$V_z$ for all \nEvtAll{} events read (grey line) and for the selected events
(blue, filled). Red dashed lines: $|V_z|<70$\,cm. 1\,cm bins.}
\hfill
\figpanel{0.48}{evt_vxvy}{$V_y$ vs $V_x$ for selected events; the red dashed circle is $V_r=2$\,cm. 0.05\,cm bins.}
\\[1ex]
\figpanel{0.48}{evt_vzdiff}{$V_z^{\rm TPC}-V_z^{\rm VPD}$ (0.1\,cm bins, log scale). The central
Gaussian core has $\sigma=\vzDiffSigma$\,cm; events without VPD vertex are outside the range.}
\hfill
\figpanel{0.48}{evt_refmult}{refMult (blue solid) and gRefMult (red dashed), 1-unit bins, log scale.
Means \meanRefMult{} and \meanGRefMult.}
\caption{Event-level distributions of the test sample.}
\end{figure}

\paragraph{Multiplicities and BTOF hits.} Per selected event the mean refMult is \meanRefMult, the mean
nBTOFMatch (\fcmd{StPicoEvent::nBTOFMatch}) is \meanNBTofMatch{} and the mean number of stored BTOF hits
is \meanNBTofHits{} (maximum \maxNBTofHits). The mean of \fcmd{btofTrayMultiplicity} is also \meanTofMult.
Figure~\ref{fig:evt_nhits_vs_nmatch} shows that the number of BTOF hits grows linearly with the number of
matched tracks, and Fig.~\ref{fig:evt_ntoftrk_vs_nhits} that the number of TOF-matched \emph{primary}
QA tracks follows it closely (correlation factor \corrTofTrkHits). On average there are \hitsPerMatched{}
BTOF hits per TOF-matched QA track; most hits are therefore \emph{not} associated with a selected track
--- they come from tracks outside the selection, secondaries, or noise, and are the
candidates for random matches.

\caveat{The hit counts and \fcmd{btofTrayMultiplicity} have identical mean and RMS in this sample, which
suggests the picoDst keeps all BTOF hits of the event. Not verified in \fcmd{StPicoDstMaker}.}

\begin{figure}[H]
\centering
\figpanel{0.48}{evt_refmult_vs_tofmult}{refMult vs \fcmd{btofTrayMultiplicity}.}
\hfill
\figpanel{0.48}{evt_refmult_vs_nmatch}{refMult vs nBTOFMatch. No pile-up cut is applied; the band
is narrow with few outliers in this sample.}
\\[1ex]
\figpanel{0.48}{evt_nhits_vs_nmatch}{Number of stored BTOF hits vs nBTOFMatch.}
\hfill
\figpanel{0.48}{evt_ntoftrk_vs_nhits}{Number of TOF-matched primary QA tracks per event
(Sec.~\ref{sec:tofdef}) vs number of stored BTOF hits.}
\\[1ex]
\figpanel{0.48}{evt_nbtofhits}{Number of stored BTOF hits per selected event, log scale.}
\caption{Event multiplicity correlations (log colour scale for 2D panels).}
\end{figure}
